\chapter{Kommentierte Quellen}

Achtundzwanzig Werke, eingeordnet: Zu jedem steht, wofür man es aufschlägt und was es leistet, das andere nicht leisten.

\lit{1}{Alon, U.: \textit{An Introduction to Systems Biology}. 2.\ Aufl., Chapman \& Hall 2019.}{Über Rückkopplungsschaltungen in Zellen --- Schalter, Schwingungen, Anpassung. Alon schreibt ungewöhnlich knapp und zeigt, wie wenige Bauteile genügen, um bestimmtes Verhalten zu erzeugen. Die Fortsetzung von Kapitel 11 und 12 auf der Ebene der Genregulation.}

\lit{2}{Anderson, R., May, R.: \textit{Infectious Diseases of Humans}. Oxford 1991.}{Das Werk, das die Epidemiemodellierung als Fach begründete. Behandelt Kapitel 7 und 8 in voller Breite --- Altersstruktur, Impfstrategien, endemische Gleichgewichte, Parasiten. Trotz des Alters unverzichtbar; die Datenbeispiele sind sorgfältig und immer noch lehrreich.}

\lit{3}{Bacaër, N.: \textit{A Short History of Mathematical Population Dynamics}. Springer 2011.}{Vierzig kurze Kapitel, je eines zu einer historischen Arbeit --- von Fibonacci über Malthus, Verhulst und Volterra bis zu Kermack und McKendrick. Jedes enthält die Rechnung im Original. Die beste Quelle für die Geschichtskästen dieses Bandes.}

\lit{4}{Bohl, E.: \textit{Mathematik in der Biologie}. 4.\ Aufl., Springer 2006.}{\textbf{Auf Deutsch und für Anfänger geeignet.} Bohl führt die nötige Mathematik parallel zu den biologischen Fragen ein und verlangt wenig Vorkenntnis. Wo dieser Band zu schnell war, ist Bohl der Einstieg.}

\lit{5}{Britton, N.: \textit{Essential Mathematical Biology}. Springer 2003.}{Eine gut aufgebaute Einführung mittleren Anspruchs, die die Themen dieses Bandes in derselben Reihenfolge behandelt --- von Populationen über Epidemien zu Erregbarkeit und Musterbildung. Mit Übungsaufgaben, allerdings ohne ausführliche Lösungen.}

\lit{6}{Derendorf, H., Gramatté, T., Schäfer, H., Staab, A.: \textit{Pharmakokinetik kompakt}. 3.\ Aufl., Wissenschaftliche Verlagsgesellschaft 2011.}{\textbf{Auf Deutsch, für Kapitel 10 die praktische Ergänzung.} Behandelt Mehrkammermodelle, nichtlineare Kinetik und die Anpassung bei Organfunktionsstörungen ausführlicher, als es hier möglich war --- mit Rechenbeispielen aus der Klinik.}

\lit{7}{Diekmann, O., Heesterbeek, H., Britton, T.: \textit{Mathematical Tools for Understanding Infectious Disease Dynamics}. Princeton 2013.}{Die theoretisch anspruchsvollste Darstellung der Epidemiemodellierung. Hier steht die allgemeine Definition von $\mathcal{R}_0$ über den größten Eigenwert der nächsten Generation --- und damit der saubere Rahmen für alles, was Kapitel 7 bis 9 nur andeuten.}

\lit{8}{Edelstein-Keshet, L.: \textit{Mathematical Models in Biology}. SIAM 2005 (Nachdruck von 1988).}{Ein Lehrbuch, das aus der Lehre entstanden ist und das merkt man. Besonders stark bei diskreten Modellen und Phasenporträts; die Kapitel zu Kapitel 3 und 6 dieses Bandes sind dort ausführlicher und mit vielen durchgerechneten Beispielen.}

\lit{9}{Ellner, S., Guckenheimer, J.: \textit{Dynamic Models in Biology}. Princeton 2006.}{Verbindet Theorie mit Rechnerarbeit --- jedes Kapitel enthält Programmieraufgaben. Für Leser, die die Modelle dieses Bandes selbst durchrechnen wollen, der praktischste Einstieg.}

\lit{10}{Ermentrout, B., Terman, D.: \textit{Mathematical Foundations of Neuroscience}. Springer 2010.}{Die Vertiefung zu Kapitel 12. Behandelt Hodgkin-Huxley vollständig, die Reduktion auf zwei Gleichungen, Bifurkationsanalyse erregbarer Systeme und gekoppelte Neuronennetze. Anspruchsvoll, aber sehr klar.}

\lit{11}{Gatenby, R., Brown, J.: \textit{Integrating evolutionary dynamics into cancer therapy}. Nature Reviews Clinical Oncology 17, 2020.}{Die Übersichtsarbeit zur adaptiven Therapie aus Kapitel 14 --- von ihrem Urheber. Lesenswert auch als Beispiel dafür, wie ein mathematisches Argument eine klinische Strategie begründet und welchen Widerstand das findet.}

\lit{12}{Goldie, J., Coldman, A.: \textit{A mathematic model for relating the drug sensitivity of tumors to their spontaneous mutation rate}. Cancer Treatment Reports 63, 1979.}{Sechs Seiten, die die Onkologie verändert haben. Die Rechnung ist die aus Satz 14.1; ihre Folgerung --- früh und kombiniert behandeln --- bestimmt bis heute den Entwurf von Therapieschemata.}

\lit{13}{Hethcote, H.: \textit{The Mathematics of Infectious Diseases}. SIAM Review 42, 2000.}{Eine Übersicht über die gesamte Modellfamilie: SIS, SIR, SEIR, mit und ohne Demographie, jeweils mit Schwellenbedingung und Gleichgewichten. Frei verfügbar und als Nachschlagewerk ausgezeichnet.}

\lit{14}{Hodgkin, A., Huxley, A.: \textit{A quantitative description of membrane current}. Journal of Physiology 117, 1952.}{Die Arbeit aus Kapitel 12. Bemerkenswert lesbar, mit vollständiger Angabe der Messungen und der Rechnung. Man sieht dem Text an, dass jede Kurve wochenlange Handrechnung gekostet hat.}

\lit{15}{Izhikevich, E.: \textit{Dynamical Systems in Neuroscience}. MIT Press 2007.}{\textbf{Das Buch, das Erregbarkeit vollständig über Bifurkationen erklärt.} Izhikevich ordnet alle bekannten Neuronen\-typen danach, welche Bifurkation ihrem Übergang zum Feuern zugrunde liegt --- ein Ord\-nungs\-prinzip, das Kapitel 12 nur streift. Reich bebildert.}

\lit{16}{Keeling, M., Rohani, P.: \textit{Modeling Infectious Diseases in Humans and Animals}. Princeton 2008.}{\textbf{Für Teil III die maßgebliche Darstellung.} Vom SIR-Modell über Altersstruktur, Raum und Netzwerke bis zu stochastischen Modellen --- mit Programmen zu jedem Abschnitt. Wer selbst Epidemiemodelle bauen will, beginnt hier.}

\lit{17}{Kermack, W., McKendrick, A.: \textit{A contribution to the mathematical theory of epidemics}. Proceedings of the Royal Society A 115, 1927.}{Die Urarbeit. Enthält Modell, Schwellensatz und Endgrößengleichung --- alles, was Kapitel 7 behandelt. Die Herleitung ist allgemeiner als die heute übliche und deshalb stellenweise schwer zu lesen; das Ergebnis ist dasselbe.}

\lit{18}{Lloyd-Smith, J. u.\,a.: \textit{Superspreading and the effect of individual variation on disease emergence}. Nature 438, 2005.}{Die Arbeit hinter Abschnitt 9.4. Zeigt an Daten mehrerer Ausbrüche, wie stark die Übertragungszahlen streuen --- und dass dieselbe Reproduktionszahl je nach Streuung völlig verschiedene Ausbruchsmuster erzeugt.}

\lit{19}{Luria, S., Delbrück, M.: \textit{Mutations of bacteria from virus sensitivity to virus resistance}. Genetics 28, 1943.}{Der Schwankungsversuch aus Kapitel 14. Methodisch ein Musterbeispiel: Eine Frage über einen unbeobachtbaren Vorgang wird durch die Verteilung beobachtbarer Ergebnisse entschieden. Frei verfügbar und erstaunlich gut lesbar.}

\lit{20}{May, R.: \textit{Simple mathematical models with very complicated dynamics}. Nature 261, 1976.}{Der Aufsatz, der das Chaos in die Biologie brachte. Fünf Seiten, mit der Aufforderung, die Gleichung selbst durchzuspielen. Mays Warnung vor überzogenen Erwartungen an Vorhersagen ist heute so aktuell wie damals.}

\lit{21}{Murray, J.: \textit{Mathematical Biology}, Band I und II. 3.\ Aufl., Springer 2002/2003.}{\textbf{Das Standardwerk des Fachs}, in zwei Bänden und rund $1500$ Seiten. Band I behandelt alles, was in diesem Band steht; Band II geht zu Musterbildung, Wellen und räumlichen Modellen über. Kein Buch zum Durchlesen, aber das erste zum Nachschlagen.}

\lit{22}{Nöbauer, W., Timischl, W.: \textit{Mathematische Modelle in der Biologie}. Vieweg 1979.}{\textbf{Auf Deutsch und knapp.} Ein älteres, aber sauber geschriebenes Buch, das die klassischen Modelle --- Wachstum, Konkurrenz, Epidemien, Genetik --- mit voll\-ständigen Rech\-nungen dar\-stellt. Für Leser, die eine deutsch\-sprachige Zweit\-darstellung suchen.}

\lit{23}{Nowak, M.: \textit{Evolutionary Dynamics}. Harvard 2006.}{Die mathematische Behandlung der Evolution: Auslese, Mutation, Spieltheorie, Wirt-Erreger-Systeme. Die Kapitel zur Virusdynamik und zur Resistenzentstehung sind die natürliche Fortsetzung von Kapitel 14.}

\lit{24}{Otto, S., Day, T.: \textit{A Biologist's Guide to Mathematical Modeling in Ecology and Evolution}. Princeton 2007.}{\textbf{Das Buch darüber, wie man ein Modell \kursiv{baut}} --- nicht wie man vorhandene löst. Führt durch den gesamten Weg von der biologischen Frage über die Annahmen zur Gleichung und zurück zur Deutung. Umfangreiche mathematische Anhänge für Leser ohne Vorkenntnisse.}

\lit{25}{Rowland, M., Tozer, T.: \textit{Clinical Pharmacokinetics and Pharmacodynamics}. 5.\ Aufl., Wolters Kluwer 2019.}{Das Standardwerk zu Kapitel 10, aus der klinischen Anwendung heraus geschrieben. Behandelt auch die Pharmakodynamik --- den Weg vom Spiegel zur Wirkung ---, der in diesem Band nur über die Hill-Kurve angedeutet wurde.}

\lit{26}{Segel, L., Edelstein-Keshet, L.: \textit{A Primer on Mathematical Models in Biology}. SIAM 2013.}{Besonders stark bei der Trennung von Zeitskalen und der Quasi-Stationarität aus Kapitel 11 --- Segel hat die zugehörige Theorie wesentlich entwickelt. Zeigt, wann die Näherung zulässig ist und wie groß ihr Fehler wird.}

\lit{27}{Strogatz, S.: \textit{Nonlinear Dynamics and Chaos}. 2.\ Aufl., Westview 2015.}{\textbf{Das beste Lehrbuch über dynamische Systeme, das es gibt.} Phasenporträts, Stabilität, Bifurkationen, Chaos --- alles, was Kapitel 3 und 6 brauchen, mit ungewöhnlich guten Erklärungen und vielen biologischen Beispielen. Auch für Leser ohne Vorkenntnisse zugänglich.}

\lit{28}{Wodarz, D., Komarova, N.: \textit{Dynamics of Cancer}. World Scientific 2014.}{Die mathematische Behandlung der Krebsentstehung und -behandlung: Mutationsakkumulation, Resistenz, Therapieschemata, Stammzellmodelle. Die Vertiefung zu den Kapiteln 13 und 14.}

\begin{merksatz}
\boxtitel{Wenn Sie nur drei Bücher kaufen}
\begin{center}
\renewcommand{\arraystretch}{1.6}
\begin{tabular}{@{}P{1.0cm}P{2.8cm}P{4.4cm}@{}}
\toprule
\textbf{[27]} & \textbf{Strogatz} & \textbf{Für die Werkzeuge}\\
\textbf{[24]} & \textbf{Otto, Day} & Für das Modellieren selbst\\
\textbf{[21]} & \textbf{Murray} & \textbf{Zum Nachschlagen}\\
\bottomrule
\end{tabular}
\end{center}

\begin{center}
\begin{minipage}{0.88\textwidth}
\kursiv{Strogatz lehrt das Denken in Phasenporträts und Bifurkationen --- danach sieht man die Kapitel 3 bis 6 dieses Bandes als eine einzige zusammenhängende Sache. Otto und Day beantworten die Frage, die kein anderes Buch ernsthaft stellt: Wie kommt man von einer biologischen Frage zu einer Gleichung? Murray ist das Nachschlagewerk, in dem alles steht.}
\end{minipage}
\end{center}

\begin{center}
\begin{minipage}{0.88\textwidth}
\kursiv{Drei kurze Ergänzungen, alle frei verfügbar: [20] May mit fünf Seiten, die man nicht vergisst, [19] Luria und Delbrück als methodisches Musterstück, und [13] Hethcote als Übersicht über die Epidemiemodelle. Auf Deutsch: [4] Bohl als Einstieg und [6] Derendorf für die Pharmakokinetik.}
\end{minipage}
\end{center}
\end{merksatz}

\bigskip
\begin{flushright}
\itshape\color{grau}
Alle Modelle sind falsch.\\
Brauchbar sind die,\\
deren Grenzen man kennt.\\[10pt]
\normalfont\color{grau}\upshape
--- der Anfang dieses Bandes\\
\itshape und sein Maßstab.
\end{flushright}
